\subsection{Whole-instruction endpoint searches}

The further physical search changes the equivalence boundary. Earlier
screens primarily require one parent binary step to match a short child
block. Here a candidate may redistribute work over a complete source
instruction before reaching the required frame. This admits a different
family of potential reductions.

The baseline fixture contains 4,959 complete source-outcome frames. All
primary registers are zero except that the tested positive decrement
register is one. These are concrete filters, unlike the unbounded
symbolic contracts used for the positive constructions. All baseline
endpoints are checked literally; 59 selected cases cover all 392 rows.
A guarded fixture uses the 274,880-pair source positivity invariant and
has 4,956 representative feasible-branch cases. It still covers every
baseline row. These representative concrete frames are not asserted to
occur in the actual blank run.

For each deleted state, a keeper is selected, all incoming edges are
redirected, and either row of the keeper may be replaced by any legal
action. The full 196-state screen examines
\[
                 196\cdot195\cdot2\cdot784=59,928,960
\]
row assignments and finds no survivor. The guarded exact version and a
relaxed-frame version each examine the 49-state engine/router subset,
with 14,982,240 assignments per version, and find none.
The relaxed version allows a different gap of at least 48 and arbitrary
valid spare blocks while retaining the source path, head, marker and
all nine primary register values.

A separate helper extension permits the changed row to belong to any
retained control. Its completed prefix consists of deleted states
\file{0,1,2,3,4,5,6,7,9,10}; it checks 29,495,648 assignments with no
survivor. The interrupted work on state 14 is excluded from that count.
Failure-trace intersections restrict which single helper rows can
possibly repair all failed tests within the same budget.

Each concrete test allows at most
$\max(800,2s+64)$ candidate transitions, where $s$ is the parent
instruction duration, and uses a 256-cell tape array. Unexpected halt,
leaving the array, or failing to reach the required endpoint within the
budget rejects the assignment. A wrong intermediate root visit alone
does not reject it. An independent set-of-cells simulator checks all
4,959 baseline endpoints and all 4,959 endpoints after a genuine two-step
head round trip is inserted before root/read-0. Deliberate halt and
infinite-motion controls are rejected for the appropriate reasons.

The search counts are assignments and overlap between lanes. A rejected
assignment may merely require a longer simulation or a different
representation. Thus the conclusion is precisely that no tested
assignment passes the stated finite endpoint filter. It is not a proof
of different blank halting behavior or a general lower bound. Complete
sources, fixtures, positive controls, limits and logs are retained under
\file{research/macro_boundary_study/}.
